% !TeX root = ../main.tex
\exercise{۲.۷}
\subsection*{صورت مسئله، داده‌ها و قرارداد محاسبه}
یک ترانزیستور \lr{nMOS} در فرایند \(\qty{0.6}{\micro\metre}\) دارای ضخامت اکسید
\(t_{\mathrm{ox}}=\qty{10}{\nano\metre}\)، معادل صد آنگستروم،
غلظت ناخالصی پذیرنده \(N_A=\qty{2e17}{\per\cubic\centi\metre}\)
و آستانه بایاس صفر \(V_{t0}=\qty{0.7}{\volt}\) است.
بدنه به زمین متصل است. اگر ولتاژ سورس از صفر به \(\qty{4}{\volt}\) برسد،
ولتاژ آستانه در دمای اتاق چقدر تغییر می‌کند
\cite[تمرین \exnum{۲.۷}، صفحه چاپی ۹۶]{weste}؟

برای سازگاری با ثابت‌ها و تقریب‌های مثال \exnum{۲.۵} کتاب، دمای اتاق را \(T=\qty{300}{\kelvin}\)
و ولتاژ حرارتی را \(v_T=kT/q\approx\qty{26}{\milli\volt}\) در نظر می‌گیریم؛
\(k\) ثابت بولتزمن و \(q\) قدرمطلق بار الکترون است.
ثابت‌های مورد استفاده عبارت‌اند از
\cite[صفحات چاپی ۶۵، ۷۹ و ۸۰]{weste}
\begin{align*}
 n_i&=\qty{1.45e10}{\per\cubic\centi\metre},
 &q&\approx\qty{1.6e-19}{\coulomb},\\
 \varepsilon_0&\approx\qty{8.85e-14}{\farad\per\centi\metre},
 &\varepsilon_{\mathrm{ox}}&=3.9\varepsilon_0,
 &\varepsilon_{\mathrm{si}}&=11.7\varepsilon_0.
\end{align*}
در این روابط، \(n_i\) غلظت ذاتی حامل‌ها، \(\varepsilon_0\) گذردهی خلأ،
\(\varepsilon_{\mathrm{ox}}\) گذردهی اکسید و \(\varepsilon_{\mathrm{si}}\) گذردهی سیلیکون است.
محاسبات میانی با ارقام بیشتر نگه داشته می‌شوند و نتیجه نهایی گرد می‌شود.

\subsection*{محاسبه پارامترهای اثر بدنه و تغییر آستانه}
از اتصال بدنه به زمین، \(V_{SB}=V_S-V_B=\qty{4}{\volt}\) به دست می‌آید.
رابطه کامل اثر بدنه همان رابطه \eqref{eq:body-threshold} است.
پتانسیل سطحی و ضریب اثر بدنه از روابط کتاب محاسبه می‌شوند
\cite[روابط \exnum{۲.۳۶} و \exnum{۲.۳۷}]{weste}\cite[اسلایدهای ۱۶ و ۱۷]{nonideal}:
\begin{equation}
 \phi_s=2v_T\ln\!\left(\frac{N_A}{n_i}\right),\qquad
 C_{\mathrm{ox}}=\frac{\varepsilon_{\mathrm{ox}}}{t_{\mathrm{ox}}},\qquad
 \gamma=\frac{\sqrt{2q\varepsilon_{\mathrm{si}}N_A}}{C_{\mathrm{ox}}}.
 \label{eq:numeric-body-parameters}
\end{equation}
با \(t_{\mathrm{ox}}=\qty{1e-6}{\centi\metre}\)، پتانسیل سطحی و ظرفیت اکسید بر واحد سطح عبارت‌اند از
\begin{align}
 \phi_s&=2(\qty{0.026}{\volt})
 \ln\!\left(\frac{2\times10^{17}}{1.45\times10^{10}}\right)
 \approx\qty{0.85486}{\volt},\label{eq:numeric-phi}\\
 C_{\mathrm{ox}}&=\frac{\qty{3.4515e-13}{\farad\per\centi\metre}}
 {\qty{1e-6}{\centi\metre}}
 =\qty{3.4515e-7}{\farad\per\square\centi\metre}.
 \label{eq:numeric-cox}
\end{align}
برای کنترل واحدها، ضریب صورت رابطه \(\gamma\) را \(B_Q\) می‌نامیم.
واحد آن از \(\unit{\farad}=\unit{\coulomb\per\volt}\) به دست می‌آید:
\begin{align}
 B_Q&\equiv\sqrt{2q\varepsilon_{\mathrm{si}}N_A}
 \approx\qty{2.57427e-7}{\coulomb\per\square\centi\metre\per\volt\tothe{0.5}},
 \label{eq:numeric-charge-coefficient}\\
 [B_Q]&=\sqrt{\unit{\coulomb}\,
 \unit{\farad\per\centi\metre}\,
 \unit{\per\cubic\centi\metre}}
 =\unit{\coulomb\per\square\centi\metre\per\volt\tothe{0.5}}.
 \label{eq:numeric-dimension}
\end{align}
پس \(B_Q\) به‌تنهایی بار بر واحد سطح نیست؛ ضرب آن در ریشه یک ولتاژ، واحد بار سطحی می‌دهد.
تقسیم آن بر \(C_{\mathrm{ox}}\)، ضریب اثر بدنه را با واحد ریشه ولتاژ تعیین می‌کند:
\begin{equation}
 \gamma=\frac{B_Q}{C_{\mathrm{ox}}}
 \approx\qty{0.74584}{\volt\tothe{0.5}}.
 \label{eq:numeric-gamma}
\end{equation}
اکنون با جایگذاری روابط \eqref{eq:numeric-phi} و \eqref{eq:numeric-gamma} در رابطه \eqref{eq:body-threshold} داریم
\begin{align}
 \Delta V_t&\equiv V_t-V_{t0}
 =\gamma\left(\sqrt{\phi_s+\qty{4}{\volt}}-\sqrt{\phi_s}\right)
 \approx\qty{0.95377}{\volt},\label{eq:numeric-threshold-shift}\\
 V_t&=\qty{0.7}{\volt}+\Delta V_t
 \approx\qty{1.65377}{\volt}.
 \label{eq:numeric-threshold-new}
\end{align}
ارقام بیشتر در این دو رابطه برای بازتولید محاسبه‌اند.
با گرد کردن متناسب با دقت تقریبی داده‌ها، پاسخ نهایی چنین است:
\begin{equation}
 \boxed{\Delta V_t\approx\qty{0.95}{\volt},\qquad V_t\approx\qty{1.65}{\volt}.}
 \label{eq:numeric-body-result}
\end{equation}
بنابراین آستانه افزایش می‌یابد.
این آستانه، ولتاژ گیت نسبت به سورس در شروع وارونگی است؛ افزایش خود ولتاژ سورس نباید دوباره به \(\Delta V_t\) اضافه شود.

\subsection*{کنترل اعتبار تقریب خطی}
خطی‌سازی رابطه اثر بدنه در اطراف \(V_{SB}=0\) می‌دهد
\cite[رابطه \exnum{۲.۳۸}]{weste}\cite[اسلاید ۱۷]{nonideal}
\begin{equation}
 k_\gamma=\left.\frac{\partial V_t}{\partial V_{SB}}\right|_{V_{SB}=0}
 =\frac{\gamma}{2\sqrt{\phi_s}}\approx0.4033,
 \qquad \Delta V_{t,\mathrm{lin}}\approx k_\gamma V_{SB}\approx\qty{1.61}{\volt}.
 \label{eq:numeric-linear-check}
\end{equation}
شرط اعتبار این تقریب، \(V_{SB}\ll\phi_s\)، در این مسئله برقرار نیست؛ نسبت این دو کمیت حدود \num{4.7} است.
تقریب خطی، تغییر آستانه را حدود ۶۹ درصد بیش‌برآورد می‌کند.
علت، مقعر بودن تابع ریشه دوم و کاهش شیب آن با افزایش بایاس است؛ خود آستانه در مدل رادیکالی به مقدار ثابت اشباع نمی‌شود.
ازاین‌رو، پاسخ این مسئله باید از رابطه کامل غیرخطی به دست آید.
